% Slajd 1 – tytuł. Nazwę zespołu, ID i adresy ustawia się w main.tex.
{
\setbeamercolor{background canvas}{bg=accNavy}
\setbeamercolor{normal text}{fg=white}
\usebeamercolor[fg]{normal text}
\begin{frame}[plain]
  \vspace{6mm}
  \begin{columns}[T,onlytextwidth]
    \begin{column}{0.64\textwidth}
      \begin{tikzpicture}
        \node[fill=white,rounded corners=2.5mm,inner sep=1mm] {\includegraphics[height=13mm]{logo}};
      \end{tikzpicture}

      \vspace{2.5mm}
      {\usebeamerfont{title}\color{white}Accessly\par}
      \vspace{0.5mm}
      {\Large\bfseries\color{white}Kraków bez barier\par}
      \vspace{2mm}
      {\color{accSignal}Twoje miasto. Twoje potrzeby. Sprawdź zanim pójdziesz.\par}
      \vspace{4mm}
      {\small\color{white}
        Przewodnik po dostępnych miejscach i~mapa barier. Każda informacja ma źródło, datę i~status weryfikacji.\par
        \vspace{3mm}
        \textbf{\TeamName}\ \textperiodcentered\ \TeamId\par
        HackYeah 2026 \textperiodcentered\ wyzwanie „Kraków bez barier” (Gmina Miejska Kraków)\par
        \vspace{2.5mm}
        {\footnotesize
        Demo: \DemoUrl\\
        Kod: \RepoUrl\\
        Film (3~min): \VideoUrl\par}}
    \end{column}
    \begin{column}{0.32\textwidth}
      \centering
      \vspace{-1mm}
      \includegraphics[height=70mm]{mobile-mapa}
    \end{column}
  \end{columns}

  \note[item]{Jedno zdanie: Accessly mówi nie „czy miejsce jest dostępne”, tylko „czy zadziała dla mnie” – i~pokazuje skąd to wiemy, od kiedy i~jak bardzo można temu ufać.}
  \note[item]{Pokazać telefon obok laptopa: ten sam prototyp działa na obu.}
  \note[item]{Nie czytać linków – są w~PDF.}
\end{frame}
}
